% Part: many-valued-logic
% Chapter: three-valued-logics
% Section: lukasiewicz

\documentclass[../../../include/open-logic-section]{subfiles}

\begin{document}

\olfileid[hi]{mvl}{inf}{luk}

\olsection{लुकासिएविच तर्कशास्त्र}

\begin{editorial}
  यह अनंत-मूल्य लुकासिएविच तर्कशास्त्र पर एक संक्षिप्त, आरंभिक खंड है।
\end{editorial}

\begin{defn}\ollabel{def:lukasiewicz} अनंत-मूल्य लुकासिएविच
तर्कशास्त्र~$\LogLuk[\infty]$ निम्नलिखित आव्यूह से परिभाषित होता है:
\begin{enumerate}
  \item $\lnot$, $\land$, $\lor$, $\lif$ वाली मानक प्रतिज्ञप्तिक भाषा
  $\Lang L_0$।
  \item सत्य-मूल्यों का समुच्चय $V_\infty$।
  \item $1$ एकमात्र निर्दिष्ट मूल्य है, अर्थात् $V^+ = \{1\}$।
  \item सत्य-फलन निम्नलिखित फलनों से दिए जाते हैं:
  \begin{align*}
    \tf{\lnot}[\LogLuk](x) & = 1 - x\\
    \tf{\land}[\LogLuk](x,y) & = \min(x,y)\\
    \tf{\lor}[\LogLuk](x,y) & = \max(x,y)\\
    \tf{\lif}[\LogLuk](x,y) & = \min(1,1-(x-y)) = \begin{cases}
      1 & \text{यदि } x \le y\\
      1-(x-y) & \text{अन्यथा।}
    \end{cases}
    \end{align*}
\end{enumerate}
$m$-मूल्य लुकासिएविच तर्कशास्त्र भी इसी प्रकार परिभाषित होता है, केवल
$V = V_m$ लिया जाता है।
\end{defn}

\begin{prop}
  \olref[thr][luk]{def:lukasiewicz} से परिभाषित तर्कशास्त्र $\LogLuk[3]$
  वही है जो \olref{def:lukasiewicz} से परिभाषित $\LogLuk[3]$ है।
\end{prop}

\begin{proof}
  इसे \olref[thr][luk]{def:lukasiewicz} में संयोजकों के लिए दी गई
  सत्य-सारणियों की तुलना \olref{def:lukasiewicz} के समीकरणों से निर्धारित
  सत्य-सारणियों के साथ करके देखा जा सकता है:
  \begin{center}
    \begin{tabular}{c|c} 
      $\tf{\lnot}$ & \\ 
      \hline  
      $1$ & $0$ \\ 
      $1/2$ & $1/2$ \\
      $0$ & $1$ 
    \end{tabular}
    \quad
    \begin{tabular}{c|ccc} 
      $\tf{\land}[\LogLuk[3]]$ & $1$ & $1/2$ & $0$ \\ 
      \hline 
      $1$ & $1$ & $1/2$ & $0$ \\ 
      $1/2$ & $1/2$ & $1/2$ & $0$\\ 
      $0$ & $0$ & $0$ & $0$ 
    \end{tabular}
    \\[2ex]
    \begin{tabular}{c|ccc} 
      $\tf{\lor}[\LogLuk[3]]$ & $1$ & $1/2$ & $0$ \\ 
      \hline 
      $1$ & $1$ & $1$ & $1$ \\ 
      $1/2$ & $1$ & $1/2$ & $1/2$ \\
      $0$ & $1$ & $1/2$ & $0$ 
    \end{tabular}
    \quad
    \begin{tabular}{c|ccc} 
      $\tf{\lif}[\LogLuk[3]]$ & $1$ & $1/2$ & $0$ \\ 
      \hline 
      $1$ & $1$ & $1/2$ & $0$ \\ 
      $1/2$ & $1$ & $1$ & $1/2$  \\ 
      $0$ & $1$ & $1$ & $1$ 
    \end{tabular}
  \end{center} 
\end{proof}

\begin{prop}\ollabel{prop:luk-infty-m}
  यदि $\Gamma \Entails[\LogLuk[\infty]] !B$, तो प्रत्येक $m \ge 2$ के लिए
  $\Gamma \Entails[\LogLuk[m]] !B$।
\end{prop}

\begin{proof}
  अभ्यास।
\end{proof}

\begin{prob}
  \olref[mvl][inf][luk]{prop:luk-infty-m} सिद्ध कीजिए।
\end{prob}

वास्तव में विलोम भी मान्य है।

अनंत-मूल्य लुकासिएविच तर्कशास्त्र सर्वाधिक लोकप्रिय फ़ज़ी तर्कशास्त्र है।
फ़ज़ी तर्कशास्त्र के साहित्य में सशर्त को प्रायः $\lnot !A \lor !B$ के रूप
में परिभाषित किया जाता है। इससे अनंत-मूल्य प्रबल क्लिनी तर्कशास्त्र प्राप्त
होगा।

\begin{prob}
  दिखाइए कि $(p \lif q) \lor (q \lif p)$, $\LogLuk[\infty]$ की
  सर्वसत्यता है।
\end{prob}

\end{document}
